\documentclass[11pt]{ltjsarticle}
\usepackage{luatexja-fontspec}
\usepackage{amsmath,amssymb}
\usepackage[margin=22mm]{geometry}
\usepackage{graphicx}
\usepackage{longtable,booktabs,array}
\usepackage{calc}
\usepackage{xcolor}
\usepackage{hyperref}
\hypersetup{colorlinks=true,linkcolor=blue,urlcolor=blue,citecolor=blue}
\makeatletter
\def\maxwidth{\ifdim\Gin@nat@width>\linewidth\linewidth\else\Gin@nat@width\fi}
\makeatother
\setkeys{Gin}{width=\maxwidth,keepaspectratio}
\setcounter{secnumdepth}{0}
\setlength{\parindent}{0pt}
\setlength{\parskip}{0.5em}
\providecommand{\tightlist}{\setlength{\itemsep}{0pt}\setlength{\parskip}{0pt}}
\providecommand{\pandocbounded}[1]{#1}
\providecommand{\real}[1]{#1}
\newcounter{none}
\begin{document}
\section{Appendix: The Logarithmic Map Separates Scale and
Quantization}\label{appendix-the-logarithmic-map-separates-scale-and-quantization}

\subsection{--- g is Flattened, ω Remains
Invariant}\label{g-is-flattened-ux3c9-remains-invariant}

\textbf{Author:} Noriaki Kihara (ORCID: 0009-0004-6753-4020)
\textbf{Date:} 18 June 2026 \textbf{DOI (this version):}
10.5281/zenodo.20742989 \textbf{Concept DOI (all versions):}
10.5281/zenodo.20742988 \textbf{Zenodo:}
https://zenodo.org/records/20742989 \textbf{Related paper:} ``On the
Logarithmic Map of Scale and the Separation of Scale and Quantization in
a High-Symmetry Limit,'' §7.3 (Concept DOI:
\href{https://doi.org/10.5281/zenodo.20740841}{10.5281/zenodo.20740841})

\begin{center}\rule{0.5\linewidth}{0.5pt}\end{center}

\subsection{Purpose}\label{purpose}

This appendix develops independently the core of the related paper, §7.3
(``The heart of the map: separation of scale (g) and quantization
(ω)''). It adds no new claim; its aim is to make explicit the most
easily misread point --- that \textbf{flattening by the logarithmic map
is not ``throwing away information'' but separating scale (g) and
quantization (ω)}.

Ordinarily, flattening in a symmetric limit is read as ``structure
vanishes / becomes trivial.'' Indeed, as stated in §7.2 of the parent
paper, the fully symmetric limit, as the price of making flattening
hold, almost exhausts the degrees of freedom on the g side. That the
physics (quantization) nonetheless does not vanish is because
quantization resided from the start not in g (the mean, scale) but in ω
(the fluctuation, area). The logarithmic map acts on g and flattens it,
and does not act on ω, leaving it invariant. This is the separation.

This appendix is not a rigorous proof but an explication of the
structure of §7.3 of the parent paper. The magnitude of the fluctuation
±1/2 is a posit (parent paper §6.2). No physical identification is made.
``Separation'' here is meant only operationally, and includes no claim
of a canonical transformation strictly preserving the symplectic form
(its rigorization is future work).

\begin{center}\rule{0.5\linewidth}{0.5pt}\end{center}

\subsection{1. The two quantities: g and
ω}\label{the-two-quantities-g-and-ux3c9}

This system has two independent real conserved quantities (parent paper,
Axioms 4 and 5; and the appendix ``The Complex Structure is a
Consequence, Not a Posit''). The two correspond to the real and
imaginary parts of the complex inner product $\langle$·,·$\rangle$=g+iω.

\textbf{g (real part, distance, scale):} the sum of squares

\[g = \sum \nu_n^2 = R^2.\]

This appears to the observer as the scale R (the size, distance, metric
of the system), the side of normalization (the Born norm). It is the
invariant of SO(2) --- the rotation preserving the sum of squares.

\textbf{ω (imaginary part, area, quantization):} the product of
fluctuations

\[\omega = \delta_1\delta_2 = \tfrac{1}{2}.\]

This is an area / symplectic form, the side carrying quantization. It is
related to SO(1,1) --- the hyperbolic, product-preserving
transformation.

By the note of the parent paper, g and ω are independent quantities,
mutually underivable. One carries scale (a continuous magnitude), the
other quantization (the discrete minimal unit).

\begin{center}\rule{0.5\linewidth}{0.5pt}\end{center}

\subsection{2. The logarithmic map acts on
g}\label{the-logarithmic-map-acts-on-g}

Introduce the fundamental quantity as the logarithm of a squared
quantity, N=log\textsubscript{2}ν (parent paper §4). This map acts on the g side.

g is the sum of squares Σν². In the fully symmetric limit (all ν\_n²
equal to a value v, parent paper §7.2),

\[g = \sum\nu_n^2 = (\text{number of axes})\cdot v,\]

so the logarithm can act on the sum:

\[\log_2 g = \log_2(\text{number of axes}) + \log_2 v.\]

That is, g (the sum of squares, scale, the curved quantity) is folded
under the logarithm into a flat additive relation N+½log\textsubscript{2}(number of
axes)=N\_R (parent paper §7.2). The source of scale is the square
(power), and the logarithm linearizes the power.

\textbf{g can be removed as a gauge.} In the fully symmetric limit g
degenerates to a single value v. Scale (size) is folded into the
logarithmic coordinate and absorbed as the choice of the flat coordinate
(a gauge). After flattening, no structure carrying quantization remains
on the g side --- because g carried scale and never carried
quantization.

\begin{center}\rule{0.5\linewidth}{0.5pt}\end{center}

\subsection{3. The logarithmic map does not act on
ω}\label{the-logarithmic-map-does-not-act-on-ux3c9}

On the other hand, ω is not crushed by the logarithmic map.

ω is the product of fluctuations δ\textsubscript{1}δ\textsubscript{2}=1/2 (or the half-bit Δ=±1/2; the
appendix ``From the Sum-Zero of Fluctuations to Uncertainty''). This is
an area --- a quantity spanned by two axes --- a different kind of
quantity from the sum of squares (the length of one axis). The logarithm
linearizes the sum of squares (distance, g), but the area (ω) is outside
its action.

Concretely, even when g degenerates to a single value v in the fully
symmetric limit, the fluctuation Δ=±1/2 does not degenerate. The
degeneration of g is ``all ν\_n² are equal,'' which is the degeneration
of the mean (scale). The fluctuation (the ±1/2 allocation around the
mean) does not vanish even when the mean becomes a single value. The
degree of freedom of which member of the conjugate pair the half-bit is
allocated to (Δ\_a+Δ\_b=0, sum-zero) is preserved independently of the
degeneration of g.

\textbf{ω remains as an invariant on the flat surface.} Before and after
the logarithmic map, ω --- the area / half-bit --- is invariant. On the
flat N-surface, ω remains in a purified form as the minimal unit of
quantization.

\begin{center}\rule{0.5\linewidth}{0.5pt}\end{center}

\subsection{4. The structure of the
separation}\label{the-structure-of-the-separation}

From the above, the logarithmic map separates the two quantities.

{\def\LTcaptype{none} % do not increment counter
\begin{longtable}[]{@{}
  >{\raggedright\arraybackslash}p{(\linewidth - 4\tabcolsep) * \real{0.3333}}
  >{\raggedright\arraybackslash}p{(\linewidth - 4\tabcolsep) * \real{0.3333}}
  >{\raggedright\arraybackslash}p{(\linewidth - 4\tabcolsep) * \real{0.3333}}@{}}
\toprule\noalign{}
\begin{minipage}[b]{\linewidth}\raggedright
Quantity
\end{minipage} & \begin{minipage}[b]{\linewidth}\raggedright
What it carries
\end{minipage} & \begin{minipage}[b]{\linewidth}\raggedright
Under the logarithmic map
\end{minipage} \\
\midrule\noalign{}
\endhead
\bottomrule\noalign{}
\endlastfoot
g (sum of squares, distance, real part) & scale (continuous magnitude) &
flattened (removable as a gauge) \\
ω (product, area, imaginary part) & quantization (discrete minimal unit)
& remains invariant \\
\end{longtable}
}

\textbf{The operational meaning of ``separation'':} ``to separate''
means that the map flattens g (scale, sum of squares) while keeping ω
(product, area) invariant. g and ω are the real and imaginary parts of
the complex inner product g+iω, and the logarithmic map acts on this
real part and does not act on the imaginary part (Figure 1).

\begin{figure}
\centering
\pandocbounded{\includesvg[keepaspectratio]{fig4_g_omega_separation.pdf}}
\caption{Figure 1: g/ω separation by the logarithmic map. The curved
scale g (sum of squares) in ν-space is flattened in N-space and
degenerates to a point (gauge). On the other hand, the quantization ω
(the half-bit ±1/2) remains invariant under the map. Flattening is not a
loss of information but a separation of g and ω.}
\end{figure}

\textbf{The substance of ``the physics does not vanish even when
flattened'':} in the fully symmetric limit, flattening makes g (scale)
degenerate to a single value. That quantization (ω) nonetheless remains
is because quantization resided from the start not in g (the mean,
scale) but in ω (the fluctuation, area). Flattening folds g but does not
fold ω. Hence, even when scale is removed, quantization remains.

This separation is a consequence of quantization (ω) being a conserved
quantity independent of, and distinct from, scale (g). Scale is pushed
into the flat surface and gauged away, and quantization is purified as
an invariant on that flat surface --- this is the work of the
logarithmic map.

\begin{center}\rule{0.5\linewidth}{0.5pt}\end{center}

\subsection{5. Summary of the structure}\label{summary-of-the-structure}

For the complex inner product g+iω, the logarithmic map N=log\textsubscript{2}ν:

\begin{itemize}
\tightlist
\item
  \textbf{acts on the real part g (sum of squares, scale) and flattens
  it in the fully symmetric limit (a gauge);}
\item
  \textbf{does not act on the imaginary part ω (product, quantization),
  and leaves it as an invariant on the flat surface.}
\end{itemize}

Flattening is not a loss of information but a separation of two
quantities. Scale (g) carries a continuous magnitude and is flattenable;
quantization (ω) carries the discrete minimal unit and is invariant
under the map. That the two are independent (the parent paper's note:
the mutual underivability of Axioms 4 and 5) is the ground for this
separation.

This is the heart of the parent paper --- even when flattened in a
high-symmetry limit, quantization does not vanish, because quantization
resides not in the side that is flattened (g, scale) but in the side
that remains invariant (ω, area).

\begin{center}\rule{0.5\linewidth}{0.5pt}\end{center}

\subsection{6. Limitations (reading
notes)}\label{limitations-reading-notes}

This appendix makes explicit the structure of §7.3 of the parent paper
and adds no new claim.

\begin{enumerate}
\def\labelenumi{\arabic{enumi}.}
\item
  \textbf{``Separation'' is meant only operationally.} It means the map
  flattens g and keeps ω invariant. It includes no claim of a canonical
  transformation strictly preserving the symplectic form. That
  rigorization (a proof that the map is a canonical transformation) is
  future work.
\item
  \textbf{±1/2 is a posit.} As in §6.2 of the parent paper, the minimal
  unit of ω, ±1/2, is a posit, not a derivation, and the separation of
  this appendix presupposes this posit.
\item
  \textbf{The difference from known frameworks is uncompared.} Whether
  the g/ω separation of this appendix differs from or agrees with known
  Kähler geometry (the structure in which a metric g and a symplectic
  form ω are tied by a complex structure) or conformal flatness has not
  been rigorously compared. As in the limitation §8.4 of the parent
  paper, whether this result remains a special case of known frameworks
  cannot be determined.
\item
  \textbf{Holds only in a limited regime.} As in §8.1 of the parent
  paper, the flattening (on the g side) holds only in the k=1, fully
  symmetric, anonymity limit. Leaving this limit, the flattening does
  not hold. The separation of this appendix is also a statement within
  this limit.
\item
  \textbf{No physical identification is made.} The physical names of g
  (scale) and ω (quantization) are anonymous labels, and this appendix
  makes no identification with a particular physical quantity.
\end{enumerate}

\begin{center}\rule{0.5\linewidth}{0.5pt}\end{center}

\subsection{Reference}\label{reference}

\begin{itemize}
\tightlist
\item
  Kihara, N. ``On the Logarithmic Map of Scale and the Separation of
  Scale and Quantization in a High-Symmetry Limit --- A Limited
  Verification Report from a Minimal Axiom System,'' §7.2, §7.3, §8.1,
  §8.4, Axioms 4 and 5. Concept DOI:
  \href{https://doi.org/10.5281/zenodo.20740841}{10.5281/zenodo.20740841}.
\end{itemize}
\end{document}
